\chapter{Future Work}

\section{Stage 2 Goal}
The next phase of DeepSceneLoc aims to move from scene-level context prediction to exact-location inference, including landmark, city, and country estimation from a single image.

\section{Planned Technical Extensions}
\subsection{Hybrid Location Reasoning}
Scene probabilities from Stage 1 will be used as priors for Stage 2 location inference, reducing search space and improving consistency of final predictions.

\subsection{Model Upgrades}
Future experiments will include EfficientNet and Vision Transformer variants to evaluate gains in global-context understanding and class boundary separation.

\subsection{Confidence Calibration}
Calibration techniques such as temperature scaling will be investigated to improve reliability of probability outputs before cross-module integration.

\subsection{Data Expansion}
Dataset diversity will be increased with additional geographic regions and seasonal variation to improve domain robustness.

\subsection{Embedding-Space Analysis}
Future work will include embedding-space diagnostics (for example, PCA and t-SNE projections) to inspect class separability and identify persistent overlap zones. This analysis will guide targeted data augmentation and hard-sample mining.

\subsection{Top-$k$ Reliability Analysis}
Because location reasoning is inherently uncertain, top-$k$ scene reliability will be evaluated alongside top-1 accuracy. This helps Stage 2 consume richer priors and reduce brittle dependence on a single label.

\section{Deployment and Usability Enhancements}
\begin{itemize}
\item Optimize inference latency for near real-time usage.
\item Enhance demo interface with top-$k$ suggestions and uncertainty indicators.
\item Add automated report generation for experiment runs.
\end{itemize}

\section{Validation Plan}
Validation in Stage 2 will include:
\begin{itemize}
\item city-level and country-level accuracy,
\item confidence-weighted correctness analysis,
\item ablation studies for stage wise contribution.
\end{itemize}

Validation will also include confidence-calibration curves, error-bucket analysis by ambiguity type, and controlled robustness tests across lighting/weather variants.

\section{Risks and Mitigation}
\begin{itemize}
\item \textbf{Risk:} limited exact-location ground truth.
\item \textbf{Mitigation:} phased benchmarking at region/city granularity.
\item \textbf{Risk:} API dependency and cost variability.
\item \textbf{Mitigation:} cached inference and controlled evaluation batches.
\item \textbf{Risk:} increased compute demand for transformer training.
\item \textbf{Mitigation:} transfer-learning schedules and selective fine-tuning.
\end{itemize}

\section{Expected Outcome}
The extended system is expected to deliver a practical hybrid geolocation framework where semantic understanding and exact-location reasoning operate as a cohesive, interpretable pipeline.

\section{Milestone-Oriented Execution Plan}
The Stage 2 execution plan is organized into measurable milestones:
\begin{itemize}
\item \textbf{Milestone 1:} complete EfficientNet-B0 and ViT benchmark runs with consistent split policy.
\item \textbf{Milestone 2:} integrate confidence calibration and top-$k$ reliability reporting.
\item \textbf{Milestone 3:} implement and validate location-reasoning API pipeline.
\item \textbf{Milestone 4:} produce end to end evaluation with city/country-level scoring.
\item \textbf{Milestone 5:} finalize dissertation-grade documentation and viva artifacts.
\end{itemize}

This milestone framing ensures that future work remains auditable and aligned with both academic deadlines and technical quality targets.

